\documentclass[11pt]{article}

\usepackage[margin=1in]{geometry}
\usepackage{amsmath,amssymb}
\usepackage{booktabs,longtable,tabularx,array}
\usepackage{graphicx}
\usepackage{float}
\usepackage{xcolor}
\usepackage{microtype}
\usepackage[hidelinks]{hyperref}
\usepackage{fancyhdr}
\setlength{\headheight}{14pt}

\definecolor{navy}{HTML}{17365D}
\definecolor{slate}{HTML}{455A64}
\newcolumntype{Y}{>{\raggedright\arraybackslash}X}
\newcolumntype{L}[1]{>{\raggedright\arraybackslash}p{#1}}
\newcommand{\dd}{\,\mathrm{d}}
\Urlmuskip=0mu plus 1mu

\pagestyle{fancy}
\fancyhf{}
\lhead{Volatility-aware portfolio mathematics}
\rhead{Vancouver, 2026}
\cfoot{\thepage}

\title{\textcolor{navy}{\textbf{Mathematical Anatomy of a Multi-Currency\\
Portfolio and Volatility-Aware Overnight Strategy}}}
\author{A visual mathematical note}
\date{Vancouver, 2026}

\begin{document}
\maketitle

\begin{abstract}
This note gives a visual mathematical introduction to a constant-weight, multi-currency ETF
portfolio with a volatility-aware SPY close-to-next-open sleeve. It develops currency conversion,
portfolio returns, maximum-likelihood estimation, realized variance, fractional-Kelly-style
sizing, performance measurement, microstructure statistics, and dependent-data bootstrap
inference. Graphs are placed beside the concepts they illuminate, including cumulative wealth,
rolling Sharpe ratios, drawdowns, volatility-conditioned decisions, sleeve contribution, and
parameter sensitivity. A final numerical example amalgamates the methods. The historical figures
describe one short research sample; they are not investment advice or evidence of future returns.
\end{abstract}

\tableofcontents

\section{Research design}\label{sec:scope}

The portfolio combines a 50\% thematic ETF sleeve with a 50\% SPY/TLT 60/40 sleeve, after moving
5\% of total capital into EUR-hedged equivalents. SPY represents 27\% of total net asset value
(NAV). That slice is either held from one close to the next open or left in cash according to a
rolling return-and-volatility rule; all other holdings remain at their strategic weights.

The saved research sample contains 730 common sessions from 1 September 2023 through
17 September 2026. Prices are expressed in CAD, non-SPY holdings are rebalanced daily, and the
overnight leg is charged a fixed two-basis-point round-trip cost. These assumptions are kept
visible because they determine what the graphs can and cannot show.

\section{A chapter-level study catalogue}\label{sec:catalogue}

The ordering in Table~\ref{tab:reading} is a study path, not a claim that every chapter is required
before reading the mathematics. It covers every method used by the strategy, while also marking
where a book goes materially beyond the present model. Chapter names and numbering were
checked against the publishers' online catalogues listed in the bibliography.

\begin{longtable}{@{}L{0.065\textwidth}L{0.27\textwidth}L{0.25\textwidth}L{0.31\textwidth}@{}}
\caption{Ordered reading catalogue mapped to the mathematical topics.}\label{tab:reading}\\
\toprule
Order & Book & Chapters & What to learn and where it appears here \\
\midrule
\endfirsthead
\toprule
Order & Book & Chapters & What to learn and where it appears here \\
\midrule
\endhead
1 & Ruppert and Matteson, \emph{Statistics and Data Analysis for Financial Engineering} \cite{ruppert2015}
& Ch.~2 ``Returns''; Ch.~5 ``Modeling Univariate Distributions''
& Simple and log returns, compounding, sample moments, Gaussian likelihood, and the difference between an estimator and a model assumption. Read first for Equations~\eqref{eq:simple-return}--\eqref{eq:mle}. \\
\addlinespace
2 & Bacon, \emph{Practical Portfolio Performance Measurement and Attribution} \cite{bacon2023}
& Ch.~3 ``The Mathematics of Portfolio Return''; Ch.~5 ``Risk''; App.~B ``Multi-Currency Attribution Methodology''
& Wealth linking, annualization, risk-adjusted performance, drawdown context, and the accounting discipline needed for Equations~\eqref{eq:fx}--\eqref{eq:sharpe}. \\
\addlinespace
3 & Ruppert and Matteson \cite{ruppert2015}
& Ch.~7 ``Multivariate Statistical Models''; Ch.~16 ``Portfolio Selection''; Ch.~17 ``The Capital Asset Pricing Model''; Ch.~19 ``Risk Management''
& Covariance, correlation, beta, portfolio weights, and risk measurement. These quantities are used descriptively; the strategy does not solve a Markowitz optimization problem or estimate CAPM expected returns. \\
\addlinespace
4 & Degiannakis and Floros, \emph{Modelling and Forecasting High Frequency Financial Data} \cite{degiannakis2015}
& Ch.~2 ``Intraday Realized Volatility Measures''; Ch.~3 ``Methods of Volatility Estimation and Forecasting''
& High-frequency log returns and realized variance, corresponding to Equation~\eqref{eq:rv}. The repository scales partial-session realized variance but does not fit a volatility forecast. \\
\addlinespace
5 & Harris, \emph{Trading and Exchanges} \cite{harris2002}
& Ch.~4 ``Orders and Order Properties''; Ch.~14 ``Bid/Ask Spreads''; Ch.~19 ``Liquidity''; Ch.~21 ``Liquidity and Transaction Cost Measurement''
& MOC, market-on-open, and market orders; full quoted spreads; liquidity; and why a constant basis-point cost is only a first approximation. These chapters contextualize Equations~\eqref{eq:spread} and \eqref{eq:strategy-return}. \\
\addlinespace
6 & Davison and Hinkley, \emph{Bootstrap Methods and their Application} \cite{davison1997}
& Ch.~2 ``The Basic Bootstraps''; Ch.~5 ``Confidence Intervals''; Ch.~8 ``Complex Dependence''
& Resampling logic, percentile intervals, and dependence-preserving schemes. This is the basis for the circular moving-block procedure in Equation~\eqref{eq:block-bootstrap}. \\
\addlinespace
7 & L\'opez de Prado, \emph{Advances in Financial Machine Learning} \cite{lopez2018}
& Ch.~2 ``Financial Data Structures''; Ch.~10 ``Bet Sizing''; Ch.~11 ``The Dangers of Backtesting''; Ch.~14 ``Backtest Statistics''; Ch.~19 ``Microstructural Features''
& Bars, discretized exposure, look-ahead and selection hazards, drawdown/Sharpe diagnostics, and spread estimators. This is best read after the basic mathematics because its principal value here is methodological caution. \\
\addlinespace
8 & Ruppert and Matteson \cite{ruppert2015}
& Ch.~12 ``Time Series Models: Basics''; Ch.~13 ``Time Series Models: Further Topics''; Ch.~14 ``GARCH Models''; Ch.~6 ``Resampling''
& Extensions not used here: serial dependence, conditional volatility, diagnostics, and more formal resampling. These chapters define a sensible next research step. \\
\bottomrule
\end{longtable}

\section{Portfolio accounting}\label{sec:portfolio}

\subsection{Currency conversion}

Let asset $i$ be quoted in currency $Q_i$ and let the account base currency be $B$. If
$X_t^{B,Q_i}$ denotes units of $Q_i$ per one unit of $B$, a native price $P_{i,t}^{Q_i}$ has
base-currency value
\begin{equation}
 P_{i,t}^{B}=\frac{P_{i,t}^{Q_i}}{X_t^{B,Q_i}}.
 \label{eq:fx}
\end{equation}
For example, a USD price is divided by the quoted CADUSD rate when the base is CAD. Equivalently,
one may use a direct multiplier $Y_t^{Q_i,B}$ measured in units of $B$ per unit of $Q_i$:
\begin{equation}
 P_{i,t}^{B}=P_{i,t}^{Q_i}Y_t^{Q_i,B}, \qquad
 Y_t^{Q_i,B}=\left(X_t^{B,Q_i}\right)^{-1}.
 \label{eq:fx-live}
\end{equation}
The two equations are identical after units are made explicit. Forward-filling an FX series onto
an asset's trading dates is an alignment convention, not an estimate of an unobserved intraday
rate. It avoids dropping a valid asset close merely because the calendars differ.

\subsection{NAV and target weights}

With cash $C_t$, position quantity $q_{i,t}$, and base-currency price $P_{i,t}^{B}$, NAV is
\begin{equation}
 V_t=C_t+\sum_{i=1}^{m}q_{i,t}P_{i,t}^{B}.
 \label{eq:nav}
\end{equation}
For target weights $w_i\geq 0$ satisfying $\sum_i w_i\leq 1$, the ideal target quantity is
\begin{equation}
 q_{i,t}^{\star}=\frac{w_iV_t}{P_{i,t}^{B}}.
 \label{eq:target-quantity}
\end{equation}
Equation~\eqref{eq:target-quantity} is the bridge between economic allocation and position size.
Any unallocated weight remains cash. The historical results below use fractional quantities so
that rounding does not obscure the portfolio mathematics.

\section{Returns, wealth, and performance}\label{sec:performance}

For a consistently adjusted price series $P_t$, the one-period simple return is
\begin{equation}
 r_t=\frac{P_t}{P_{t-1}}-1,
 \label{eq:simple-return}
\end{equation}
and the log return is
\begin{equation}
 \ell_t=\log\!\left(\frac{P_t}{P_{t-1}}\right)=\log(1+r_t).
 \label{eq:log-return}
\end{equation}
Simple returns aggregate naturally across simultaneously held assets; log returns add through
time and are used by the intraday realized-variance calculation. Mixing adjusted closes with raw
opens would create artificial overnight returns, so the overnight module explicitly requires
consistent adjustment.

Starting from wealth $V_0$, compounded wealth after $T$ returns is
\begin{equation}
 V_T=V_0\prod_{t=1}^{T}(1+r_t).
 \label{eq:wealth}
\end{equation}
Total return and compound annual growth rate over $Y$ calendar years are
\begin{equation}
 R_{\mathrm{tot}}=\frac{V_T}{V_0}-1,
 \qquad
 \operatorname{CAGR}=\left(1+R_{\mathrm{tot}}\right)^{1/Y}-1.
 \label{eq:cagr}
\end{equation}

\begin{figure}[H]
\centering
\includegraphics[width=0.94\textwidth]{../graphics_paper/figures/fig1_equity.pdf}
\caption{Cumulative value of the full portfolio and principal comparators. A logarithmic vertical
scale makes equal proportional changes occupy equal distances.}
\label{fig:equity}
\end{figure}

Using $A=252$ trading sessions per year and a zero risk-free rate, the annualized
volatility and Sharpe ratio as
\begin{equation}
 \widehat{\sigma}_{\mathrm{ann}}=\widehat{\sigma}_{r}\sqrt{A},
 \qquad
 \widehat{S}=\frac{\bar r}{\widehat{\sigma}_{r}}\sqrt{A}.
 \label{eq:sharpe}
\end{equation}
For the rolling graph, $\bar r$ and $\widehat\sigma_r$ are recomputed over the latest 126 sessions
using the sample standard deviation. The rolling estimate is highly path-dependent and can change
quickly when an extreme observation enters or leaves the window.

\begin{figure}[H]
\centering
\includegraphics[width=0.94\textwidth]{../graphics_paper/figures/fig9_rolling_sharpe.pdf}
\caption{Rolling 126-session annualized Sharpe ratios. The common rise and fall across strategies
shows how strongly the statistic depends on the prevailing market window.}
\label{fig:rolling-sharpe}
\end{figure}

For a paired comparison at session $t$, define
\begin{equation}
 \Delta \widehat S_t=widehat S_t^{\mathrm{vol\text{-}aware}}
 -\widehat S_t^{\mathrm{SPY}}.
 \label{eq:rolling-sharpe-difference}
\end{equation}
A positive value means that the volatility-aware portfolio had the larger Sharpe ratio over that
particular trailing window. Because adjacent 126-session windows overlap almost completely, the
histogram is a descriptive view of the path rather than a sample of independent observations.

\begin{figure}[H]
\centering
\includegraphics[width=0.82\textwidth]{../graphics_paper/figures/fig10_sharpe_difference.pdf}
\caption{Distribution of the paired rolling Sharpe difference in
Equation~\eqref{eq:rolling-sharpe-difference}. Values to the right of zero favour the
volatility-aware portfolio within the corresponding historical window.}
\label{fig:sharpe-difference}
\end{figure}

Let $M_t=\max_{0\leq s\leq t}V_s$ be the running high-water mark. Drawdown and maximum drawdown
are
\begin{equation}
 D_t=\frac{V_t}{M_t}-1,
 \qquad
 D_{\max}=\min_{0\leq t\leq T}D_t.
 \label{eq:drawdown}
\end{equation}
These summaries describe one realized path. They do not quantify parameter uncertainty,
multiple testing, tail severity beyond the observed sample, or the probability of future loss.

\begin{figure}[H]
\centering
\includegraphics[width=0.94\textwidth]{../graphics_paper/figures/fig2_drawdown.pdf}
\caption{Drawdown paths from each strategy's own previous peak.}
\label{fig:drawdown}
\end{figure}

\section{The overnight signal}\label{sec:signal}

\subsection{Return definitions and information timing}

Let $O_t$ and $C_t$ denote the adjusted open and close of session $t$. The overnight and daily
simple returns are
\begin{equation}
 r_t^{\mathrm{ON}}=\frac{O_t}{C_{t-1}}-1,
 \qquad
 r_t^{\mathrm{D}}=\frac{C_t}{C_{t-1}}-1.
 \label{eq:overnight-daily}
\end{equation}
A decision submitted near the close of day $t-1$ cannot use $O_t$ or $C_t$. The backtest's slice
at sale day $t$ estimates overnight moments from observations ending at $O_{t-1}$ and daily
moments from closes ending at $C_{t-2}$. This explicit lag is the main protection against
look-ahead bias in the signal calculation.

\subsection{Local Gaussian maximum likelihood}

For the most recent $n$ overnight observations $x_1,\ldots,x_n$, assume locally and only as a
working approximation that $x_j\overset{\mathrm{iid}}{\sim}\mathcal N(\mu,\sigma^2)$. The
Gaussian log-likelihood, up to an additive constant, is
\begin{equation}
 \mathcal L(\mu,\sigma^2)
 =-\frac{n}{2}\log(\sigma^2)-\frac{1}{2\sigma^2}
 \sum_{j=1}^{n}(x_j-\mu)^2.
 \label{eq:likelihood}
\end{equation}
Maximization yields the estimators used here,
\begin{equation}
 \widehat{\mu}=\frac{1}{n}\sum_{j=1}^{n}x_j,
 \qquad
 \widehat{\sigma}_{\mathrm{ON}}^2
 =\frac{1}{n}\sum_{j=1}^{n}(x_j-\widehat{\mu})^2.
 \label{eq:mle}
\end{equation}
The divisor $n$ is the maximum-likelihood convention, not the unbiased sample-variance
convention. The corresponding estimated standard error of the mean is
\begin{equation}
 \widehat{\operatorname{se}}(\widehat{\mu})
 =\frac{\widehat{\sigma}_{\mathrm{ON}}}{\sqrt n}.
 \label{eq:se}
\end{equation}
With all-in round-trip cost $c$ as a decimal return and user-selected $z\geq 0$, the entry
condition is
\begin{equation}
 \widehat{\mu}-c>z\,
 \frac{\widehat{\sigma}_{\mathrm{ON}}}{\sqrt n}.
 \label{eq:entry-filter}
\end{equation}
At the default $z=0$, Equation~\eqref{eq:entry-filter} is merely a positive estimated-net-edge
filter; it is not a significance test and provides no type-I error guarantee.

\subsection{Realized intraday variance}

Suppose completed intraday bars have closes $p_1,\ldots,p_K$ and $p_0$ is the session opening
price. Form bar log returns $u_k=\log(p_k/p_{k-1})$. If $H$ is the full regular-session
duration and $h$ is elapsed covered time, the partial-session realized variance is scaled as
\begin{equation}
 \widehat v_{t}^{\mathrm{ID}}
 =\frac{H}{h}\sum_{k=1}^{K}u_k^2.
 \label{eq:rv}
\end{equation}
Only bars completed by the decision time are allowed. Coverage below 90\% returns a missing value,
and the trading rule skips the day rather than imputing risk. The linear $H/h$ scaling assumes
variance accrues proportionally through time; the familiar U-shaped intraday volatility pattern
makes this a practical approximation, not an identity.

\begin{figure}[H]
\centering
\includegraphics[width=0.93\textwidth]{../graphics_paper/figures/fig4_decisions.pdf}
\caption{The daily volatility reading, the resulting hold/skip decision, and the rolling share of
days held. Missing intraday readings are omitted and therefore never generate a trade.}
\label{fig:decisions}
\end{figure}

\begin{figure}[H]
\centering
\includegraphics[width=0.94\textwidth]{../graphics_paper/figures/fig5_vol_vs_overnight.pdf}
\caption{Intraday volatility versus the following overnight return, shown as individual outcomes
and volatility quintiles. The wide dispersion cautions against reading the filter as a return
forecast.}
\label{fig:vol-next-return}
\end{figure}

The sizing denominator is deliberately conservative:
\begin{equation}
 \widehat v_t=\max\left\{
 \widehat{\sigma}_{\mathrm{ON}}^2,
 \widehat{\sigma}_{\mathrm{D}}^2,
 \widehat v_t^{\mathrm{ID}}
 \right\}.
 \label{eq:conservative-variance}
\end{equation}
It is important to interpret Equation~\eqref{eq:conservative-variance} correctly. It is not an
MLE for a common variance. It is a hand-built upper envelope designed to reduce exposure when any
one of three volatility readings is large.

\subsection{Fractional-Kelly-style sizing and binary allocation}

For a small risky return $R$ and fraction $f$ of wealth invested, a second-order approximation to
expected log growth is
\begin{equation}
 \mathbb E[\log(1+fR)]
 \approx f\,\mathbb E[R]-\frac{f^2}{2}\mathbb E[R^2].
 \label{eq:log-growth}
\end{equation}
When the mean is small relative to volatility,
$\mathbb E[R^2]\approx\operatorname{Var}(R)$, giving the familiar approximate optimizer
\begin{equation}
 f^{\star}\approx\frac{\mathbb E[R]}{\operatorname{Var}(R)}.
 \label{eq:kelly-approx}
\end{equation}
The strategy uses this shape as a sizing heuristic, with fraction $\lambda=0.25$, cost $c$, variance
from Equation~\eqref{eq:conservative-variance}, and cap $f_{\max}$:
\begin{equation}
 f_t^{\mathrm{raw}}=
 \begin{cases}
 \min\!\left(f_{\max},\lambda\dfrac{\widehat\mu-c}{\widehat v_t}\right),
 & \text{if Equation~\eqref{eq:entry-filter} holds and }\widehat v_t>10^{-12},\\[6pt]
 0,&\text{otherwise.}
 \end{cases}
 \label{eq:raw-weight}
\end{equation}
This is ``Kelly-style,'' not an exact Kelly solution: return distributions are not established,
parameters are estimated, the variance is an upper envelope, and the allocation rule is then
discretized. In binary mode with threshold $f_{\min}$,
\begin{equation}
 f_t^{\mathrm{bin}}=f_{\max}
 \mathbf 1\!\left\{f_t^{\mathrm{raw}}\geq f_{\min}\right\}.
 \label{eq:binary-weight}
\end{equation}
Equation~\eqref{eq:binary-weight} applies only to the preassigned SPY sleeve; it does not lever the
entire account by $f_t^{\mathrm{raw}}$.

If the sleeve trades, its net return ending at the next open is
\begin{equation}
 r_{t}^{\mathrm{slice}}
 =f_t^{\mathrm{bin}}\left(r_t^{\mathrm{ON}}-c\right).
 \label{eq:strategy-return}
\end{equation}
If $r_t^{\mathrm{rest}}$ is the whole-account return contributed by all non-SPY target weights,
including idle cash, and $w_{\mathrm{SPY}}=0.27$, total return is
\begin{equation}
 r_t^{\mathrm{total}}
 =r_t^{\mathrm{rest}}+w_{\mathrm{SPY}}r_t^{\mathrm{slice}}.
 \label{eq:combined-return}
\end{equation}

\begin{figure}[H]
\centering
\includegraphics[width=0.92\textwidth]{../graphics_paper/figures/fig3_slice_contribution.pdf}
\caption{Cumulative contribution of the alternative overnight rules relative to leaving the SPY
slice in cash. This isolates the sleeve from the rest of the strategic portfolio.}
\label{fig:slice-contribution}
\end{figure}

\section{Microstructure screen and dependent-data uncertainty}\label{sec:screen}

For synchronized best bid $b_{i,t}$ and ask $a_{i,t}$, the full quoted spread in basis points is
\begin{equation}
 s_{i,t}=10^4\frac{a_{i,t}-b_{i,t}}{(a_{i,t}+b_{i,t})/2}.
 \label{eq:spread}
\end{equation}
The screen compares a candidate $i$ with SPY using paired overnight returns. Its descriptive
market beta and correlation are
\begin{equation}
 \widehat\beta_i=
 \frac{\widehat{\operatorname{Cov}}(r_i,r_{\mathrm{SPY}})}
 {\widehat{\operatorname{Var}}(r_{\mathrm{SPY}})},
 \qquad
 \widehat\rho_i=
 \frac{\widehat{\operatorname{Cov}}(r_i,r_{\mathrm{SPY}})}
 {\widehat\sigma_i\widehat\sigma_{\mathrm{SPY}}}.
 \label{eq:beta-corr}
\end{equation}
Let $\delta_d=\bar s_{\mathrm{SPY},d}-\bar s_{i,d}$ be the daily paired spread saving. Because
adjacent days may be dependent, the screen resamples circular blocks of length $L=5$. With a random
start $J_k$ on a circle of $n$ days, bootstrap replication $b$ uses the index sequence
\begin{equation}
 I_b=\left((J_1+0{:}L-1)\bmod n,\ldots,
 (J_K+0{:}L-1)\bmod n\right)_{1:n}
 \label{eq:block-bootstrap}
\end{equation}
and computes $\bar\delta_b^{\star}=n^{-1}\sum_{j\in I_b}\delta_j$. The 2.5\% and 97.5\%
empirical quantiles of 3,000 replications form a percentile interval. A candidate qualifies only
if beta and correlation thresholds pass and the lower interval endpoint exceeds the required raw
spread saving. The interval is exploratory: this construction does not correct for screening many
candidates, optimize block length, or account for quote-selection error.

\section{An amalgamation example}\label{sec:example}

This example uses round numbers to join the portfolio, signal, volatility, and return equations.
It is illustrative rather than a reconstruction of any account or a trading recommendation.

Assume NAV is CAD~250,000 and the target SPY sleeve is 27\%. Its capital budget is
\begin{equation}
 B_{\mathrm{SPY}}=0.27\times 250{,}000=\text{CAD }67{,}500.
 \label{eq:example-budget}
\end{equation}
Suppose SPY is USD~690 and the direct rate is 1.43 CAD per USD. By
Equation~\eqref{eq:fx-live}, the base-currency price is
\begin{equation}
 P_{\mathrm{SPY}}^{\mathrm{CAD}}=690\times1.43
 =\text{CAD }986.70.
 \label{eq:example-fx}
\end{equation}

Now take a 60-session window with estimated mean overnight return
$\widehat\mu=0.00050$, all-in cost $c=0.00020$, overnight volatility 0.50\%, daily volatility
1.20\%, and scaled intraday volatility 0.90\%. The conservative variance is
\begin{equation}
 \widehat v=\max\{0.005^2,0.012^2,0.009^2\}=0.000144.
 \label{eq:example-variance}
\end{equation}
With $z=0$, the estimated net edge $0.00030$ passes
Equation~\eqref{eq:entry-filter}. Quarter sizing gives
\begin{equation}
 f^{\mathrm{raw}}=\min\!\left(1,
 0.25\frac{0.00050-0.00020}{0.000144}\right)
 =0.5208.
 \label{eq:example-raw}
\end{equation}
Because $0.5208\geq f_{\min}=0.5$, binary mode sets $f^{\mathrm{bin}}=1$ and allocates the whole
sleeve. If the intraday volatility had instead been 1.40\%, then
$\widehat v=0.014^2=0.000196$ and $f^{\mathrm{raw}}=0.3827$; the same estimated mean would be
skipped. This is the precise sense in which the rule is volatility-aware.

Assume the next overnight return is 0.40\% and the non-SPY holdings contribute 0.15\% to
whole-portfolio return. Equations~\eqref{eq:strategy-return} and
\eqref{eq:combined-return} give
\begin{equation}
 r^{\mathrm{slice}}=1(0.0040-0.0002)=0.0038,
 \qquad
 r^{\mathrm{total}}=0.0015+0.27(0.0038)=0.002526.
 \label{eq:example-return}
\end{equation}
Thus the illustrative whole portfolio gains 25.26 basis points under the stated accounting model.

\section{Evidence and limitations}\label{sec:limitations}

Figures~\ref{fig:equity}--\ref{fig:drawdown} show a transformation of exposure and risk, not
evidence of dominance. In this sample, continuously holding the blended portfolio ends above the
volatility-aware variant, while SPY has both the largest terminal value and the deepest drawdown.
Figure~\ref{fig:rolling-sharpe} also shows that risk-adjusted rankings vary substantially through
time and depend on the chosen 126-session window.

\begin{figure}[H]
\centering
\includegraphics[width=0.92\textwidth]{../graphics_paper/figures/fig7_sensitivity.pdf}
\caption{Sensitivity of the volatility-aware strategy to the binary allocation threshold.}
\label{fig:sensitivity}
\end{figure}

The principal limitations are:

\begin{itemize}
 \item The sample is short, begins in a rising-market regime, and was viewed while the rule was
 designed; it is not an untouched out-of-sample test.
 \item The accounting omits slippage, commissions, taxes, financing, market impact, and auction
 uncertainty.
 \item The local IID Gaussian assumption is a working estimator, not a proven model of data
 generation. Financial returns may be dependent, heteroskedastic, skewed, and heavy-tailed.
 \item Sensitivity plots reveal parameter dependence but are not a substitute for out-of-sample
 validation.
 \item The overnight return is computed in USD while the total portfolio is reported in CAD, so
 overnight FX movement in the SPY leg is omitted.
 \item The partial-day variance scaler ignores deterministic intraday seasonality and
 microstructure noise. A volatility model or seasonality adjustment would be a distinct method,
 not a cosmetic change.
 \item The screen's percentile interval is not corrected for multiple candidate searches. Beta
 similarity does not make two ETFs economically interchangeable.
 \item Sharpe ratios use a zero risk-free rate; different risk-free assumptions and annualization
 windows can change their values and rankings.
\end{itemize}

A disciplined next study would freeze parameters, add explicit turnover and cost accounting,
validate the FX leg, reserve a genuinely untouched test period, and report uncertainty across
regimes.

\section*{Conclusion}

The strategy forms a coherent mathematical chain: consistent prices produce returns; returns
produce local mean and volatility estimates; a cost-adjusted, variance-scaled rule controls a
preallocated sleeve; and compounded returns produce familiar performance summaries. The graphs
make the consequences visible. In particular, the conservative maximum variance, binary
threshold, fixed cost, and historical sample are design choices rather than consequences of
financial theory. Labelling them explicitly is the most important mathematical safeguard.

\begin{thebibliography}{9}

\bibitem{ruppert2015}
David Ruppert and David S. Matteson.
\newblock \emph{Statistics and Data Analysis for Financial Engineering: with R Examples}, 2nd ed.
\newblock Springer, 2015.
\newblock Publisher catalogue and chapter list:
\url{https://link.springer.com/book/10.1007/978-1-4939-2614-5}.

\bibitem{bacon2023}
Carl R. Bacon.
\newblock \emph{Practical Portfolio Performance Measurement and Attribution}, 3rd ed.
\newblock Wiley, 2023.
\newblock Publisher catalogue and chapter list:
\url{https://onlinelibrary.wiley.com/doi/book/10.1002/9781119831976}.

\bibitem{degiannakis2015}
Stavros Degiannakis and Christos Floros.
\newblock \emph{Modelling and Forecasting High Frequency Financial Data}.
\newblock Palgrave Macmillan, 2015.
\newblock Publisher catalogue and chapter list:
\url{https://link.springer.com/book/10.1057/9781137396495}.

\bibitem{harris2002}
Larry Harris.
\newblock \emph{Trading and Exchanges: Market Microstructure for Practitioners}.
\newblock Oxford University Press, 2002.
\newblock Publisher catalogue and chapter list:
\url{https://academic.oup.com/book/52292}.

\bibitem{davison1997}
A. C. Davison and D. V. Hinkley.
\newblock \emph{Bootstrap Methods and their Application}.
\newblock Cambridge University Press, 1997.
\newblock Publisher catalogue and chapter list:
\url{https://doi.org/10.1017/CBO9780511802843}.

\bibitem{lopez2018}
Marcos L\'opez de Prado.
\newblock \emph{Advances in Financial Machine Learning}.
\newblock Wiley, 2018.
\newblock Publisher table of contents:
\url{https://catalogimages.wiley.com/images/db/pdf/9781119482086.toc.pdf}.

\bibitem{kelly1956}
John L. Kelly, Jr.
\newblock A new interpretation of information rate.
\newblock \emph{Bell System Technical Journal}, 35(4):917--926, 1956.
\newblock \url{https://doi.org/10.1002/j.1538-7305.1956.tb03809.x}.

\end{thebibliography}

\vfill
\begin{center}
\footnotesize\color{slate}
\rule{0.78\textwidth}{0.4pt}\\[0.6em]
\textbf{AI disclaimer.} This paper was tool-assisted through AI; the ideas, strategy choices,
and responsibility for their use are human. Vancouver, 2026.
\end{center}

\end{document}
